\section{Prikry forcing and change of cofinality}
\label{sec:p06-prikry-forcing-cofinality}

Classical Prikry forcing separates cofinality from cardinality: it makes a
ground measurable cardinal singular of cofinality \(\omega\), adds no
bounded subsets of that cardinal, and collapses no cardinal. The finite
stem grows under the forcing order, while a second, direct-extension
order only shrinks the measure-one tail. Normality connects the two
orders through a finite-sequence homogeneity construction. The basic
source routes are \cite[\S1.1, Lemmas~1.1--1.13]{GitikPrikryAuthorDraft}
and \cite[\S9, version of 24 August 2026]{Karagila2026ForcingNotes}.
The author draft, the published Handbook chapter and Prikry's original
1970 presentation \cite{Prikry1970OriginalMetadata} retain their separate
identities. The native genericity route is also developed in
\cite[Theorem~2.2 and Lemmas~2.5--2.6]{Mathias1973PrikryNative}; its
normality and finite-homogeneity arguments are expanded below.

\subsection{The ground measure and the two orders}

\begin{hbconvention}{def:p06-prikry-ground}{Internal normality and an externally supplied generic}
Use the transitive-ground framework of \cref{def:p06-ftr-ground}.
Thus \(M\) is a transitive set model of \(\mathsf{ZFC}\), and
\(\kappa,U\in M\) satisfy, inside \(M\), that \(\kappa\) is an
uncountable cardinal and \(U\) is a nonprincipal \(\kappa\)-complete
normal ultrafilter on \(\kappa\). All ground subsets, colorings,
intersections and completeness assertions below are computed in \(M\).
In particular \(U\) measures \(\mathcal P(\kappa)^M\), not every
ambient subset of \(\kappa\). When discussing an extension, supply an
ambient \(M\)-generic filter \(G\) for the poset defined below.
External countability of \(M\) suffices to produce such a filter; it is
not a premise of the internal combinatorial proofs. Neither transitivity
nor \(\operatorname{Con}(\mathsf{ZFC})\) supplies a normal measure or
a transitive ground with one. Cardinal and cofinality comparisons mean
those of \(M\) and \(M[G]\), not ambient cardinalities of these sets.
\end{hbconvention}

\begin{hbdefinition}{def:p06-prikry-normal-measure}{Normal measure and diagonal intersection}
Nonprincipality and \(\kappa\)-completeness have the meaning of
\cref{def:p06-measurable-cardinal}. Normality means: if \(A\in U\) and
\(f:A\to\kappa\) is regressive at nonzero points, then
\(\{\alpha\in A:f(\alpha)=\xi\}\in U\) for some \(\xi<\kappa\).
For \(A_\xi\in U\), define
\[
 \mathop{\triangle}_{\xi<\kappa}A_\xi
 =\{\beta<\kappa:(\forall\xi<\beta)\ \beta\in A_\xi\}.
\]
The quantifier is over \(\xi<\beta\), not over every \(\xi<\kappa\).
Mathias's native Definition~1.3 instead takes finite-stem tail thinning
as the normality clause. The next lemma proves the equivalence of that
clause, diagonal closure and the regressive-map clause for this measure.
\end{hbdefinition}

\begin{hblemma}{thm:p06-prikry-measure-tools}{proved=annex:p06.prikry.measure-tools}{Uniformity, inaccessibility, and finite-stem diagonalization}
Every \(A\in U\) has cardinality \(\kappa\) and is unbounded.
The ground cardinal \(\kappa\) is regular and a strong limit.
The displayed diagonal intersection belongs to \(U\). More generally,
given \(B_t\in U\) for all finite increasing stems \(t\), there is
\(D\in U\) such that
\[
 \beta\in D,\quad \max(t)<\beta
       \quad\Longrightarrow\quad \beta\in B_t.
\]
For the empty stem one can require \(D\subseteq B_{\varnothing}\)
separately. The regressive, diagonal-closure and finite-stem-thinning
normality clauses are equivalent for the standing nonprincipal
\(\kappa\)-complete ultrafilter. This supplies the simultaneous tail
thinning used later.
\end{hblemma}

\begin{hblemma}{thm:p06-prikry-homogeneity}{proved=annex:p06.prikry.homogeneity}{Homogeneity at every finite arity}
If \(A\in U\), \(\gamma<\kappa\), and
\(c:[A]^{<\omega}\to\gamma\), there is \(H\in U\), \(H\subseteq A\),
such that \(c\) is constant on \([H]^n\) for each \(n<\omega\).
The constants for different arities need not be the same. The proof
derives the finite-arity property from normality rather than importing
an unproved Rowbottom theorem.
\end{hblemma}

\begin{hbdefinition}{def:p06-prikry-poset}{Finite stem and measure-one upper part}
A condition in \(P_U\) is \(p=(s_p,A_p)\), where \(s_p\) is a finite
strictly increasing sequence in \(\kappa\), \(A_p\in U\), and every
element of \(A_p\) is larger than every element of \(s_p\).
For the empty stem this last requirement is vacuous. Put
\[
\begin{aligned}
 q\leq p\quad&\Longleftrightarrow\quad
 s_q\text{ end-extends }s_p,\quad A_q\subseteq A_p,\quad
 \operatorname{ran}(s_q)\setminus\operatorname{ran}(s_p)\subseteq A_p,\\
 q\leq^*p\quad&\Longleftrightarrow\quad s_q=s_p,\quad A_q\subseteq A_p.
\end{aligned}
\]
Smaller conditions are stronger in both orders. The full order \(\leq\)
is the forcing order; \(\leq^*\) is only the auxiliary direct order.
Write \(A\setminus(\alpha+1)\) for the tail above \(\alpha\).
The greatest condition is \((\varnothing,\kappa)\).
Gitik's author draft writes stronger conditions with \(\geq\) and
\(\geq^*\); this local construction uses the opposite order glyphs to
match the existing recursive-forcing convention. Its \(q\leq p\)
means exactly Gitik's \(q\geq p\), with the same stem and tail clauses.
Mathias's native definition also places stronger conditions on the left
of \(\leq\), but uses finite subsets and a separate stem-height
notation. These source presentations are compared explicitly below.
\end{hbdefinition}

\begin{hbdefinition}{def:p06-prikry-gitik-orders}{Gitik's finite-set presentation and native order conventions}
In Gitik's presentation, \(P\) consists of pairs \(\langle p,A\rangle\),
where \(p\) is a finite subset of \(\kappa\), \(A\in U\), and
\(\min(A)>\max(p)\). For an empty first coordinate, use the vacuous
tail convention of \cref{def:p06-prikry-poset}. The native orders are
\[
\begin{aligned}
 \langle p,A\rangle\geq\langle q,B\rangle
 \quad&\Longleftrightarrow\quad
 p\text{ end-extends }q,\quad A\subseteq B,\quad p\setminus q\subseteq B,\\
 \langle p,A\rangle\geq^*\langle q,B\rangle
 \quad&\Longleftrightarrow\quad p=q,\quad A\subseteq B.
\end{aligned}
\]
For nonempty \(q\), end extension means
\(p\cap(\max(q)+1)=q\); every finite \(p\) end-extends the empty set.
Here \(\geq\) and \(\geq^*\) mean stronger and directly stronger,
respectively. Gitik also writes their converse weak orders as \(\leq\)
and \(\leq^*\). Thus, within this native finite-set presentation,
\(r\leq t\) means \(t\geq r\). These native weak orders are not the
smaller-is-stronger orders on sequence-stem conditions in
\cref{def:p06-prikry-poset}. The different operand presentations and order
orientations are part of the definitions
\cite[Definitions~1.1--1.3]{GitikPrikryAuthorDraft}.
\end{hbdefinition}

\begin{hblemma}{thm:p06-prikry-order-bridge}{proved=annex:p06.prikry.order-bridge}{The finite-set and sequence-stem forcing correspondence}
Sending \(\langle p,A\rangle\) to \((s_p,A)\), where \(s_p\) is the
unique increasing enumeration of \(p\), is a bijection of conditions.
It preserves the stronger-than and directly-stronger-than relations when
each is read in its own convention. It carries compatible conditions,
dense requirements and generic filters in both directions. With the
corresponding recursive translation of names, the forcing relations and
their generic evaluations agree.
\end{hblemma}

\begin{hbtheorem}{thm:p06-prikry-orders}{proved=annex:p06.prikry.orders}{Closure only in the direct order, and the chain condition}
These are partial orders with \(\leq^*\subseteq\leq\).
Every direct descending sequence of length less than \(\kappa\) has
a direct lower bound. Conditions with the same stem are compatible,
so \(P_U\) has the \(\kappa^+\)-chain condition.
The full order is not countably closed: a descending sequence of
strictly lengthening finite stems has no finite-stem lower bound.
\end{hbtheorem}

\subsection{The generic sequence and the Prikry lemma}

\begin{hbtheorem}{thm:p06-prikry-sequence}{proved=annex:p06.prikry.sequence}{The new cofinal sequence and reconstruction of the filter}
The union \(C_G=\bigcup\{s_p:p\in G\}\) is a strictly increasing
\(\omega\)-sequence cofinal in \(\kappa\), and it is eventually in
every \(A\in U\). It determines the entire generic filter:
\[
 G=\{(s,A)\in P_U:s\text{ is an initial segment of }C_G,
            \ C_G\text{ after }s\text{ lies wholly in }A\}.
\]
Consequently \(M[G]=M[C_G]\), \(C_G\notin M\), and
\(\operatorname{cf}^{M[G]}(\kappa)=\omega\).
\end{hbtheorem}

\begin{hbtheorem}{thm:p06-prikry-lemma}{proved=annex:p06.prikry.lemma}{Decision without lengthening the stem}
For every \(p\in P_U\), fixed forcing formula \(\varphi\), and tuple
of ground names, some \(q\leq^*p\) decides \(\varphi\) in the
\emph{full-order} recursive forcing relation.
The proof colors finite stem additions by their possible decisions,
diagonalizes the selected decision tails and homogenizes every finite
arity. Its final step is the already established density closure of
\cref{thm:p06-ftr-density}, not a definition of forcing by generic truth.
\end{hbtheorem}

\begin{hblemma}{thm:p06-prikry-strong-dense}{proved=annex:p06.prikry.strong-dense}{A uniform finite level entering a dense open set}
For a dense downward-open \(D\subseteq P_U\) and \(p=(s,A)\), there
are \(H\in U\), \(H\subseteq A\), and \(n<\omega\) such that
for every \(m\geq n\) and \(t\in[H]^m\),
\[
 (s\mathbin{\frown}t,\ H\setminus(\max(t)+1))\in D.
\]
At \(m=0\), interpret this as \((s,H)\in D\). This stronger dense-set
statement has its own construction, not an appeal to the decision lemma.
\end{hblemma}

\begin{hbtheorem}{thm:p06-prikry-no-bounded}{proved=annex:p06.prikry.no-bounded}{No new bounded subsets, including no new reals}
For every ground \(\lambda<\kappa\),
\[
 \mathcal P(\lambda)^{M[G]}=\mathcal P(\lambda)^M.
\]
More precisely, if \(p\forces\dot X\subseteq\check\lambda\), some
direct extension of \(p\) forces \(\dot X=\check Y\) for \(Y\in M\).
The proof makes fewer than \(\kappa\) successive direct decisions and
uses an actual intersection lower bound. The unbounded sequence \(C_G\)
is therefore new without adding a real.
\end{hbtheorem}

\subsection{Preservation and the exact change}

\begin{hbtheorem}{thm:p06-prikry-preservation}{proved=annex:p06.prikry.preservation}{All cardinals and the cofinality classification}
Every ground cardinal remains a cardinal in \(M[G]\), and these are all
the cardinals of the extension. For each ground limit ordinal \(\delta\),
\[
 \operatorname{cf}^{M[G]}(\delta)=
 \begin{cases}
 \omega,&\operatorname{cf}^{M}(\delta)=\kappa,\\
 \operatorname{cf}^{M}(\delta),&\operatorname{cf}^{M}(\delta)\ne\kappa.
 \end{cases}
\]
The proof for cardinals below \(\kappa\) uses bounded graph codes; the
proof at \(\kappa\) additionally uses that it is a ground limit cardinal.
Above \(\kappa\), maximal antichains bound the possible values of a
function name. Cofinality \(\kappa\), unlike cardinality \(\kappa\),
is deliberately not preserved.
\end{hbtheorem}

\begin{hbcorollary}{thm:p06-prikry-measure-loss}{proved=annex:p06.prikry.measure-loss}{A singular strong limit, and loss of ground-measure completeness}
The extension still regards \(\kappa\) as an uncountable strong-limit
cardinal, but it is no longer measurable because it is singular.
The old set \(U\) is not countably complete for all families in
\(M[G]\): the ground tails above the successive Prikry points form
a new countable family of members of \(U\) with empty intersection.
This does not contradict its completeness for families belonging to \(M\).
\end{hbcorollary}

\begin{hbtheorem}{thm:p06-prikry-mathias}{proved=annex:p06.prikry.mathias}{The eventual-measure-one genericity criterion}
Suppose an ambient strictly increasing \(\omega\)-sequence \(C\subseteq\kappa\)
is eventually in every \(A\in U\). Write \(G(C)\) for the set of
conditions described in \cref{thm:p06-prikry-sequence} with \(C\) in place
of \(C_G\). Then \(G(C)\) is \(M\)-generic for \(P_U\).
Eventual containment forces cofinality because every ground tail belongs
to \(U\). Merely being a cofinal \(\omega\)-sequence is insufficient.
This is the Mathias criterion in the native 1973 paper and Gitik's
source route, at the present transitive-set-ground hypotheses. Mathias
also discusses proper-class grounds as theorem/proof schemata; no
proper-class extension construction is asserted here.
\end{hbtheorem}

\begin{hbexample}{def:p06-prikry-order-failure}{The next point is not directly decidable}
Choose distinct \(\alpha,\beta\in A_p\). Full extensions appending
\(\alpha\) and \(\beta\) force different next Prikry points. A direct
extension of \(p\) still has a measure-one, unbounded tail, so it cannot
force the next point to be any one specified ordinal. The Prikry lemma
decides each fixed proposition, such as equality to a specified ordinal;
it does not decide an arbitrary \(\kappa\)-valued name by one direct
extension, or allow a fusion of \(\kappa\) decisions. This distinguishes
the Prikry property from closure of the full forcing order.
\end{hbexample}

\begin{proofboundary}{def:p06-prikry-source-boundary}{Classical normal-measure forcing is one source component}
The proof does not assume \(\mathsf{GCH}\), but it does assume the stated
normal measure in the ground. It does not derive measurability from
\(\operatorname{Con}(\mathsf{ZFC})\), or prove a bare relative-consistency
implication by postulating a countable transitive ground. It does not
replace Prikry forcing by a collapse or Cohen forcing: no cardinal is
collapsed and no real is added, while \(\operatorname{cf}(\kappa)\) changes.
Tree/non-normal, one-point, supercompact, extender, Magidor, Radin and
iteration variants require their own conditions and preservation proofs.
\end{proofboundary}
